\section{Aplicação prática em Python}
\label{s.python_application}

\begin{frame}{Aplicação prática em Python}
	\justify 
	\begin{itemize}
		\item<1> O principal objetivo deste \textbf{tutorial} é construir uma \textbf{aplicação prática} em \textbf{Python};
		\\~\\
		\item<2> Aplicação será um \textbf{pacote/biblioteca}, hospedado no GitHub e compartilhado no repositório oficial PyPI;
		\\~\\
		\item<3> Pacote/biblioteca será composto por arquivos básicos para \textbf{instalação}, \textbf{dependências}, \textbf{documentação} e \textbf{testes unitários};
		\\~\\
		\item<4> Por fim, serão apresentados alguns exemplos de como \textbf{melhorar} a \textbf{divulgação} e \textbf{atrair} o uso da \textbf{comunidade}.
	\end{itemize}
\end{frame}

\subsection{Criação de repositório}
\label{ss.repository_creation}

\begin{frame}{Criação de repositório}
	\begin{figure}
		\centering
		\includegraphics[scale=0.125]{figs/video_play.png}
	\end{figure}
\end{frame}

\subsection{Implementação}
\label{ss.implementation}

\begin{frame}{Implementação}
\end{frame}

\subsection{Submissão ao PyPI}
\label{ss.pypi_submission}

\begin{frame}{Submissão ao PyPI}
\end{frame}

\begin{frame}{}
	\begin{figure}
		\centering
		\includegraphics[scale=0.125]{figs/video_play.png}
	\end{figure}
\end{frame}

\subsection{Locais de divulgação}
\label{ss.places_sharing}

\begin{frame}{Locais de divulgação}
\end{frame}